\subsection{作用律片段}

设 \((G, e, \theta, m)\) 是李群芽，X 是类为 \(C^r\) 的流形。

定义 6. 类为 \(\mathbf{C}^r\) 的 \(\mathbf{G}\) 在 \(\mathbf{X}\) 上的左作用律片段，是一个映射 \(\psi\)，定义在 \(\Omega\) 上；该开集位于 \(\mathbf{G} \times \mathbf{X}\) 中，包含 \(\{e\} \times \mathbf{X}\)，且该映射取值于 \(\mathbf{X}\)，并具有下列性质：

\begin{enumerate}
\def\labelenumi{(\roman{enumi})}
\item
  \(\psi\) 是类为 \(C^{r}\) 的；
\item
  对所有 \(x \in X\)，\(\psi(e, x) = x\)；
\item
  存在一个邻域 \(\Omega_1\)，它是 \(\{e\} \times \{e\} \times \mathbf{X}\) 在 \(\mathbf{G} \times \mathbf{G} \times \mathbf{X}\) 中的邻域，使得对于 \((g, g', x) \in \Omega_1\)，\(m(g, g')\)、\(\psi(m(g, g'), x)\) 和 \(\psi(g, \psi(g', x))\) 都有定义，并且 \(\psi(g, \psi(g', x)) = \psi(m(g, g'), x)\)。
\end{enumerate}

类为 \(\mathbf{C}^r\) 的右作用律片段类似定义。

我们通常用 \(gx\) 代替 \(\psi(g, x)\)。

设 \(G'\) 是 G 的李子群芽，\(X'\) 是 X 的子流形。假设集合 \(\Omega'\)，即由 \((g,x)\in \Omega \cap (\mathrm{G}'\times \mathrm{X}')\) 中满足 \(\psi (g,x)\in \mathrm{X}'\) 的元素组成的集合，在 \(\mathrm{G}'\times \mathrm{X}'\) 中是开集（若 \(\mathbf{X}'\) 在 \(\mathbf{X}\) 中是开集，则这一条件总能满足）。则 \(\psi |\Omega '\) 是类为 \(\mathrm{C}^r\) 的左作用律片段，作用于 \(\mathrm{G}'\) 和 \(\mathrm{X}'\)，称为由 \(\psi\) 限制到 \(\mathrm{G}'\) 和 \(\mathrm{X}'\) 所得到的作用律片段。

命题 23. 设 \((\mathbf{G}, e, \theta, m)\) 是李群芽，\(\mathbf{X}\) 是类为 \(\mathbf{C}^r\) 的流形，\(x_0\) 是 \(\mathbf{X}\) 中的一点，\(\Omega\) 是 \((e, x_0)\) 在 \(\mathbf{G} \times \mathbf{X}\) 中的开邻域，\(\psi\) 是从 \(\Omega\) 到 \(\mathbf{X}\) 的映射，具有下列性质：

\begin{enumerate}
\def\labelenumi{(\roman{enumi})}
\item
  \(\psi\) 是类为 \(C^{r}\) 的；
\item
  \(\psi(e, x)\) 在 x 充分接近 \(x_{0}\) 时等于 x；
\item
  \(\psi(m(g, g'), x) = \psi(g, \psi(g', x))\)，其中 \((g, g', x)\) 充分接近 \((e, e, x_0)\)。
\end{enumerate}

则存在开邻域 \(\mathbf{X}'\)，它是 \(x_0\) 在 \(\mathbf{X}\) 中的开邻域，以及开子集 \(\Omega'\)，它是 \(\Omega \cap (\mathrm{G} \times \mathrm{X}')\) 的开子集，使得 \(\psi|\Omega'\) 是类为 \(\mathbf{C}^r\) 的左作用律片段，作用于 \(\mathbf{G}\) 和 \(\mathbf{X}'\)。

存在开邻域 \(\mathbf{X}'\)，它是 \(x_0\) 在 \(\mathbf{X}\) 中的开邻域，以及开邻域 \(\mathbf{G}'\)，它是 \(e\) 在 \(\mathbf{G}\) 中的开邻域，使得 \(\psi(e, x) = x\) 对所有 \(x \in \mathbf{X}'\) 成立，并且

\[
\psi (g, \psi (g ^ {\prime}, x)) = \psi (m (g, g ^ {\prime}), x)
\]

对于 \((g, g', x) \in \mathrm{G}' \times \mathrm{G}' \times \mathrm{X}'\)。令 \(\Omega'\) 为由 \((g, x) \in \Omega \cap (\mathrm{G}' \times \mathrm{X}')\) 中满足 \(\psi(g, x) \in \mathrm{X}'\) 的元素所成的集合。于是 \(\Omega'\) 在 \(\mathrm{G} \times \mathrm{X}'\) 中是开集，并且 \(\mathrm{X}'\)、\(\Omega'\) 具有命题中的性质。

引理 3. 设 \(\mathbf{X}\) 是正规空间，\((\mathbf{X}_i)_{i\in I}\) 是 \(\mathbf{X}\) 的局部有限开覆盖。对所有 \((i,j)\in I\times I\) 以及所有 \(x\in \mathbf{X}_i\cap \mathbf{X}_j\)，令 \(\mathrm{V}_{ij}(x)\) 为一个 \(x\) 的邻域，包含于 \(\mathbf{X}_i\cap \mathbf{X}_j\)。于是可以为每个 \(x\in X\) 关联 \(\mathrm{V}(x)\)，它是 \(x\) 的一个邻域，使下列条件成立：

\begin{enumerate}
\def\labelenumi{(\alph{enumi})}
\item
  若 \(x \in \mathbf{X}_i \cap \mathbf{X}_j\)，则 \(\mathrm{V}(x) \subset \mathrm{V}_{ij}(x)\)；
\item
  若 \(\mathrm{V}(x)\) 与 \(\mathrm{V}(y)\) 相交，则存在 \(i\in \mathbf{I}\)，使 \(\mathrm{V}(x)\cup \mathrm{V}(y)\subset \mathrm{X}_i\)
\end{enumerate}

存在开覆盖 \(\langle \mathrm{X}_i\rangle_{i\in I}\)，它覆盖 \(\mathbf{X}\)，使得 \(\overline{\mathbf{X}}_i' \subset \mathbf{X}_i\) 对所有 \(i \in I\) 成立（《General Topology》，第 IX 章，第 § 4 节，定理 3）。令 \(x \in \mathbf{X}\)。令 \(\mathrm{V}_1(x)\) 为 \(\mathrm{V}_{ij}(x)\) 与 \(\mathbf{X}_k'\) 的交集，其中这些集合包含 \(x\)；这是 \(x\) 的开邻域。令 \(\mathrm{V}_2(x)\) 为 \(x\) 的一个邻域，包含于 \(\mathrm{V}_1(x)\) 且只与有限个 \(\mathbf{X}_i\) 相交。于是 \(\mathrm{V}_2(x)\) 只与有限个 \(\overline{\mathbf{X}}_i'\) 相交，因此下式中的集合

\[
\mathrm{V} (x) = \mathrm{V} _ {2} (x) \cap \bigcap_ {i \in I, x \notin \overline {{\mathrm{X}}} _ {i} ^ {\prime}} (\mathrm{X} - \overline {{\mathrm{X}}} _ {i} ^ {\prime})
\]

是 \(x\) 的邻域。若 \(x \in \mathrm{X}_1 \cap \mathrm{X}_j\)，则 \(\mathrm{V}_1(x) \subset \mathrm{X}_1 \cap \mathrm{X}_j\)，从而 \(\mathrm{V}(x) \subset \mathrm{X}_1 \cap \mathrm{X}_j\)。令 \(x, y\) 属于 \(\mathrm{X}\)，并假设 \(\mathrm{V}(x)\) 与 \(\mathrm{V}(y)\) 相交。存在 \(i \in I\)，使 \(x \in \mathrm{X}_i'\)。于是 \(\mathrm{V}_1(x) \subset \mathrm{X}_i'\)，从而 \(\mathrm{V}(x) \subset \mathrm{X}_i'\)，并且 \(\mathrm{V}(y) \cap \overline{\mathrm{X}_i'} \neq 0\)。于是 \(y \in \overline{\mathrm{X}_i'}\)，根据 \(\mathrm{V}(y)\) 的定义，从而 \(y \in \mathrm{X}_i\) 且 \(\mathrm{V}(y) \subset \mathrm{X}_i\)。因此 \(\mathrm{X}_i\) 包含 \(\mathrm{V}(x)\) 和 \(\mathrm{V}(y)\)。

命题 24. 设 G 是李群芽，X 是类为 \(C^{r}\) 的流形，且 (X \(_{i}\) ) \(_{i \in I}\) 是 X 的局部有限开覆盖。对所有 i ∈ I，令 \(\psi_{i}\) 为 G 上类为 \(C^{r}\) 的左作用律片段，作用于 X \(_{i}\)。假设 X 的底层拓扑空间是正规空间，并且对于所有 (i, j) ∈ I × I 以及所有 x ∈ X \(_{i}\) ∩ X \(_{j}\)，\(\psi_{i}\) 与 \(\psi_{j}\) 在 (e, x) 的一个邻域上重合。存在左作用律片段 \(\psi\)，它是 G 在 X 上的类为 \(C^{r}\) 的左作用律片段，使得对所有 i ∈ I 以及所有 x ∈ X \(_{i}\)，\(\psi_{i}\) 与 \(\psi\) 在 (e, x) 的一个邻域上重合。

对所有 \((i,j)\in \mathrm{I}\times \mathrm{I}\) 以及所有 \(x\in \mathrm{X}_i\cap \mathrm{X}_j\)，选取 \(\mathrm{V}_{ij}(x)\)，它是 \(x\) 的开邻域且包含于 \(\mathrm{X}_i\cap \mathrm{X}_j\)，使 \(\psi_i\) 和 \(\psi_j\) 在 \(\{e\} \times \mathrm{V}_{ij}(x)\) 在 \(\mathrm{G}\times \mathrm{X}\) 中的一个邻域上有定义且相等。对所有 \(x\in \mathrm{X}\)，选取 \(\mathrm{V}(x)\)，它是 \(x\) 在 \(\mathrm{X}\) 中的开邻域，使引理 3 的条件 (a) 和 (b) 得到满足。令 \(\mathbf{I}_x\) 为由 \(i\in \mathrm{I}\) 且 \(x\in \mathrm{X}_i\) 的 i 所成的集合。这是有限集。令 \(\mathrm{U}_x\) 为由 \((g,y)\in \mathrm{G}\times \mathrm{V}(x)\) 所成的集合，其中 \(\psi_i\)（\(i\in \mathrm{I}_x\)）在 \((g,y)\) 的一个邻域上有定义且彼此重合。于是 \(\mathrm{U}_x\) 是开集且 \((e,x)\in \mathrm{U}_x\)。对于 \(\psi_i\)（\(i\in \mathrm{I}_x\)），它们在 \(\mathrm{U}_x\) 上都有相同的限制。令 \(x,y\) 属于 \(\mathrm{X}\)。若 \(\mathrm{U}_x\) 与 \(\mathrm{U}_y\) 相交，则 \(\mathrm{V}(x)\) 与 \(\mathrm{V}(y)\) 相交，因而存在 \(i\in \mathrm{I}\) 使

\[
\mathrm{V} (x) \cup \mathrm{V} (y) \subset \mathrm{X} _ {i}.
\]

于是 \(i \in \mathrm{I}_x, i \in \mathrm{I}_y, \psi_i | \mathrm{U}_x = \psi_x, \psi_i | \mathrm{U}_y = \psi_y\)，从而

\[
\psi_ {x} | (\mathrm{U} _ {x} \cap \mathrm{U} _ {y}) = \psi_ {y} | (\mathrm{U} _ {x} \cap \mathrm{U} _ {y}).
\]

因此，\(\psi_x\) 定义了映射 \(\psi\)，它从 \(\mathbf{U} = \bigcup_{x\in \mathbf{X}}\mathbf{U}_x\) 到 \(\mathbf{X}\)，且 \(\mathbf{U}\) 是 \(\{e\} \times \mathbf{X}\) 在 \(\mathbf{G}\times \mathbf{X}\) 中的开邻域。显然，\(\psi\) 是类为 \(C^r\) 的，且 \(\psi (e,x) = x\) 对所有 \(x\in \mathbf{X}\) 成立。对于所有 \(i\in \mathbf{I}\) 及所有 \(x\in \mathrm{X}_i\)，\(\psi\) 与 \(\psi_x\) 及 \(\psi_i\) 在 \((e,x)\) 的一个邻域上重合，于是 \(\psi\) 满足定义 6 的条件 (iii)。

